\documentclass[12pt]{article}
\usepackage[utf8]{inputenc}
\usepackage[T1]{fontenc}
\usepackage[margin=1in]{geometry}
\usepackage{setspace}
\usepackage{xcolor}
\usepackage[colorlinks=true, urlcolor=blue!60!black, linkcolor=blue!60!black]{hyperref}
\setlength{\parindent}{0em}
\setlength{\parskip}{0.6em}
\renewcommand{\baselinestretch}{1.0}
\begin{document}
\begin{center}
{\LARGE \textit{Williamson v. Lee Optical of Oklahoma, Inc.}, 348 U.S. 483 (1955)\par}
\vspace{0.5em}
{\large Case Brief --- Two-Page Edition\par}
\end{center}

\section{Memory Jogger}
Economic regulations survive due process and equal protection challenges whenever a court can conceive of a plausible rationale, even if the law is needless, wasteful, or reforms one step at a time (\href{https://www.courtlistener.com/opinion/105285/williamson-v-lee-optical-of-oklahoma-inc/}{348 U.S. 483}).

\section{Detailed Case Facts}
A 1953 Oklahoma statute barred opticians from fitting lenses into new frames or duplicating lenses without a prescription from an ophthalmologist or optometrist, exempted sellers of ready-to-wear glasses, banned advertising of eyeglasses and frames, and barred retail stores from renting space to eye-care providers. Lee Optical sued state officials.

\section{Procedural History}
A three-judge District Court struck down the prescription requirement, the ready-to-wear exemption, the frame-advertising ban, and the retail space ban, while upholding the lens-advertising ban. Both sides appealed directly to the Supreme Court.

\section{Judicial Votes}
8--0. Douglas wrote for a unanimous Court; there were no separate opinions, and Harlan did not participate.

\section{Holding}
Every challenged provision is constitutional. A business regulation is valid if there is an evil at hand and the legislature might rationally think its measure a reasonable correction; economic classifications fail only if they are ``invidious discrimination.''

\section{Analysis of Opinions}
Douglas conceded that the prescription rule ``may exact a needless, wasteful requirement in many cases,'' but the legislature might have thought frequent eye examinations necessary to detect disease, and balancing is its job. Classification admits of ``no doctrinaire definition,'' so the legislature could address one phase of the problem at a time. Frame advertising might lure customers into buying lenses, and barring eye care from retail stores could free the profession from ``all taints of commercialism.'' The days of \textit{Lochner}-style invalidation were gone.

\section{Examples: Future Applications}
\textbf{Same-Side:} requiring casket retailers to hold funeral director licenses is upheld on a hypothesized consumer-protection rationale. \textbf{Opposite-Side:} a racial bar on optician licenses triggers strict scrutiny. \textbf{Fence-Sitter:} 1,500 hours of cosmetology school for hair braiders, with no evidence of health risk.

\section{Critique}
Libertarians such as Randy Barnett and the Institute for Justice call \textit{Lee Optical} a rubber stamp for special-interest cartels. Defenders see it as a necessary bulwark against reviving \textit{Lochner}, leaving economic tradeoffs to elected legislatures.

\section{Key Quotations}
\begin{quote}
``But the law need not be in every respect logically consistent with its aims to be constitutional. It is enough that there is an evil at hand for correction, and that it might be thought that the particular legislative measure was a rational way to correct it.''
\end{quote}
\noindent --- Justice Douglas, Opinion of the Court, 348 U.S. at 487--488.

\end{document}
